% !TEX program = xelatex
\documentclass[11pt, a4paper]{article}

\usepackage{fontspec}
\setmainfont[
  Path=/usr/share/texmf/fonts/opentype/public/tex-gyre/,
  UprightFont=texgyrepagella-regular.otf,
  BoldFont=texgyrepagella-bold.otf,
  ItalicFont=texgyrepagella-italic.otf,
  BoldItalicFont=texgyrepagella-bolditalic.otf
]{TeX Gyre Pagella}

\usepackage{amsmath, amssymb, amsthm}
\usepackage[margin=1.25in]{geometry}
\usepackage{setspace}
\onehalfspacing
\usepackage{parskip}
\setlength{\parindent}{0pt}
\setlength{\parskip}{0.85em}
\usepackage[hidelinks]{hyperref}
\usepackage{titlesec}
\titleformat{\section}{\large\bfseries}{}{0em}{}[\vspace{0.2em}\hrule\vspace{0.4em}]
\titleformat{\subsection}{\normalsize\bfseries\itshape}{}{0em}{}
\usepackage{epigraph}
\setlength{\epigraphwidth}{0.62\textwidth}
\renewcommand{\epigraphflush}{center}

% Arabic support

\theoremstyle{definition}
\newtheorem{definition}{Definition}
\newtheorem{theorem}{Theorem}
\newtheorem{proposition}{Proposition}
\newtheorem{corollary}{Corollary}
\theoremstyle{remark}
\newtheorem{remark}{Remark}

\usepackage{titling}
\setlength{\droptitle}{-2em}

\title{\textbf{Language as Constraint Field}\\
{\large Grammar, Repair, and Admissibility\\
in Semantic Manifolds}}
\author{Flyxion \\ \small Independent Researcher}
\date{2026}

\begin{document}

\maketitle

\begin{abstract}
Grammar does not generate utterances. Grammar maintains the admissibility
structure that makes utterances receivable. This essay develops a formal
theory of language as a constraint field: a dynamical system in which
utterances are operators on a discourse state space, grammaticality is
a property of trajectories rather than objects, and linguistic competence
is the ability to navigate admissibility manifolds rather than to apply
production rules. The theory inverts the generative direction. Where
Chomskyan grammar asks what rules produce well-formed sentences, the
constraint-field approach asks what constraints make utterances receivable
--- a shift from production to admissibility as the organizing primitive.
The framework is developed through three case studies. Arabic root-and-pattern
morphology, reinterpreted as a repair ecology in which the triliteral root
functions as a semantic attractor and the morphological patterns are repair
operators maintaining coherence under contextual deformation. Spanish verbal
agreement, reinterpreted as a multi-field coherence condition that native
speakers satisfy by inhabiting the field rather than consulting rules. And
conversation itself, reinterpreted as alternating CLIO projection and
semantic repair, in which each turn maintains the discourse manifold by
committing a trajectory that remains within the admissibility region. The
mathematical framework connects directly to the Relativistic Scalar-Vector
Plenum field theory (providing the field-dynamic substrate), CLIO projection
dynamics (providing the discourse-reduction operator), and the Semantic
Stability Theorem (providing the branch-coherence condition for grammatical
constructions). Language death is repair collapse. Language change is manifold
deformation. The gesture-before-symbol thesis establishes that meaning exists
prior to its linguistic encoding, and language is the repair layer that
preserves access to shared meaning against the attenuation introduced by
distance, time, and individual variation.
\end{abstract}

\newpage

\epigraph{\textit{A grammatical utterance is not one that was successfully
generated. It is one that remains inside the discourse manifold.}}{}

\section{The Wrong Question}

The dominant tradition in formal linguistics --- descending through Chomsky's
generative grammar, the minimalist program, and their derivatives --- takes
as its central question: what rules generate well-formed sentences? The
grammar is a production system. Grammaticality is a property of outputs.
A sentence is grammatical if and only if it is generable by the grammar.

This question has been enormously productive. It forced the formalization of
syntactic structure, revealed the existence of deep universals in language
acquisition, and connected linguistics to computational theory. But it has
a deep orientation problem: it begins with production rather than reception.
It asks what mechanisms a speaker uses to generate sentences, not what
conditions a discourse context imposes on receivable utterances.

The repair-equilibrium framework developed in companion papers suggests that
this orientation is the source of several persistent difficulties in
generative theory. If language is primarily a repair infrastructure --- a
system for maintaining shared semantic access against the attenuation
introduced by distance, time, and individual variation --- then the
generative question is not the right first question. The right first question
is: what constraints make an utterance receivable?

This shift from production to admissibility is not merely terminological. It
changes what the primitive objects of the theory are, what counts as an
explanation, and what the empirical targets of the theory should be. A
generative grammar is a mechanism for producing sentences; its primitive
objects are symbols, rules, and derivations. A constraint-field grammar is
a dynamical system for maintaining discourse coherence; its primitive objects
are fields, operators, and admissibility regions. The difference is the
difference between asking how a sentence is built and asking why a sentence
is receivable at this moment in this conversation by this community.

\section{The Discourse Manifold}

\begin{definition}[Discourse state space]
Let $X$ be a \textit{discourse state space}: the space of all possible
discourse configurations, where a configuration specifies the propositional
content established so far, the discourse referents in play, the illocutionary
commitments of each participant, the contextually salient alternatives, and
the pragmatic norms active in the exchange. A discourse state $x \in X$ is
a complete specification of the conversational situation at a moment.
\end{definition}

\begin{definition}[Admissibility manifold]
The \textit{discourse admissibility manifold} $\mathcal{A}(x) \subseteq X$
at discourse state $x$ is the set of states reachable from $x$ by utterances
that maintain coherence: that are semantically interpretable in the context,
that preserve the referential and illocutionary structure established by
prior discourse, and that satisfy the pragmatic norms active in the exchange.
\end{definition}

\begin{definition}[Utterance as operator]
An utterance $u$ acts as an operator
\[
u: X \to X, \quad x \mapsto u(x),
\]
transforming the discourse from state $x$ to a new state $u(x)$. The utterance
is \textit{grammatically admissible} in state $x$ if and only if $u(x) \in
\mathcal{A}(x)$: the resulting discourse state is within the admissibility
manifold of the current state.
\end{definition}

Grammaticality is therefore a property of trajectories through discourse
state space, not a property of sentences considered in isolation. The sentence
``the dog bit the man'' is admissible in some discourse states and inadmissible
in others --- not because its syntactic structure varies but because the
admissibility manifold varies with context. ``Colorless green ideas sleep
furiously'' is syntactically well-formed but semantically inadmissible in
most discourse states because no standard interpretation maps it to a
trajectory within the manifold of coherent discourse.

\begin{theorem}[Dissolution of the Grammaticality/Acceptability Gap]
In the constraint-field framework, the generative distinction between
grammaticality and acceptability is not a distinction between two properties
of sentences but a distinction between two discourse contexts in which the
same sentence is evaluated. More precisely: a sentence $s$ that is
``grammatical but unacceptable'' in generative terms is, in constraint-field
terms, a sentence that is admissible under the abstract grammatical constraint
field $\mathcal{A}_G$ but inadmissible under the concrete discourse
admissibility manifold $\mathcal{A}(x)$ at the current state $x$. A sentence
that is ``acceptable but ungrammatical'' is admissible under $\mathcal{A}(x)$
but exits the boundaries of $\mathcal{A}_G$. The apparent gap arises because
the generative framework holds the grammatical constraint field fixed and
treats discourse context as noise; the constraint-field framework treats
$\mathcal{A}(x)$ as the fundamental object and $\mathcal{A}_G$ as the
degenerate case in which contextual constraints are minimized.
\end{theorem}

\begin{proof}[Worked examples]
\textit{Example 1: Grammatical but unacceptable.}
``The horse raced past the barn fell'' is syntactically well-formed (a reduced
relative clause construction, fully derivable in any generative grammar).
But in most discourse states $x$, $u(x) \notin \mathcal{A}(x)$: the
interpretation that the utterance targets is not accessible without backtracking,
because the initial parse trajectory ``the horse raced'' commits to a main-
clause reading before the reduced relative can be recovered. The utterance is
inadmissible not because it violates a rule but because the discourse state at
the moment of parse failure has been driven outside $\mathcal{A}(x)$ by the
garden-path effect. Grammatical competence (knowledge of $\mathcal{A}_G$) does
not prevent the boundary crossing; the trajectory through discourse state space
crosses the boundary before the full structure is assembled.

\textit{Example 2: Acceptable but ungrammatical.}
``Him and me went to the store'' is widely stigmatized as ungrammatical
(accusative pronouns in subject position violate Case assignment in standard
generative accounts). But in informal registers, the utterance $u(x)$ lands
inside $\mathcal{A}(x)$: it is fully interpretable, referentially clear, and
pragmatically appropriate. What has happened is that the admissibility manifold
$\mathcal{A}(x)$ for informal spoken English has expanded to include accusative
subject pronouns, while the prescriptive constraint field $\mathcal{A}_G$
(constructed from formal written norms) has not. The ``ungrammaticality'' is
a property of the mismatch between two admissibility manifolds, not a property
of the sentence.

\textit{Example 3: Performance errors.}
Slips of the tongue, agreement attraction, and garden-path errors are all
cases where the utterance operator $u$ follows the locally strongest constraint
gradient rather than the globally correct one --- where $u(x)$ exits
$\mathcal{A}(x)$ because the gradient of the constraint field was locally
distorted by discourse salience. These are not distinct from grammaticality
failures; they are the same phenomenon (trajectory exiting the admissibility
manifold) occurring at different timescales and for different reasons. The
generative competence/performance distinction is, in the constraint-field
framework, the distinction between $\mathcal{A}_G$ (the idealized, context-
free admissibility region) and $\mathcal{A}(x)$ (the actual, context-sensitive
admissibility manifold at each discourse state). Both are admissibility
manifolds; they differ in what constraints they encode.
\end{proof}

\subsection{Syntax as Constraint Propagation}

Syntactic structure, in the constraint-field framework, is the record of how
admissibility constraints propagate through an utterance. A subject-verb
agreement constraint is not a rule that is checked at a derivational step;
it is a constraint that propagates from the subject constituent to the verb
constituent through the admissibility field, requiring that the verb occupy
a region of the field consistent with the person and number established by
the subject.

Violation of agreement is not a failure to apply a rule. It is a trajectory
that exits the admissibility manifold: the verb's position in the discourse
state space is inconsistent with the constraints established by the subject.
The discourse state $u(x)$ is outside $\mathcal{A}(x)$.

This reframing immediately explains several features of syntactic processing
that are puzzling on the generative account. Agreement attraction errors ---
in which speakers produce agreement with a local noun rather than the
grammatical head --- occur precisely when the local noun is more salient in
the current discourse state, temporarily distorting the admissibility field.
The error is not a failure of rule application but a trajectory that followed
the locally strongest constraint gradient rather than the grammatically
correct one. The gradient was distorted by the discourse context.

\section{CLIO as Discourse Reduction}

The CLIO operator --- constraint-based limitation of interpretive output ---
provides the formal model of how utterances reduce the discourse state space.

At every conversational turn, the speaker faces a space $X$ of possible
meanings and must commit to an utterance that maps the current discourse
state to a state within the admissibility manifold. This is a projection
operation:
\[
\pi_u: X \to \mathcal{A}(x), \quad \pi_u(X) = u(x).
\]
The utterance selects, from the enormous space of possible meanings, a
trajectory that lands in the admissible region. The listener performs the
inverse:
\[
\pi_u^{-1}: \mathcal{A}(x) \to X',
\]
reconstructing from the utterance a trajectory in semantic space --- not the
unique trajectory the speaker intended, but one constrained to be consistent
with the utterance and the discourse context.

Conversation is therefore alternating projection and reconstruction under
shared constraint. The speaker projects from semantic space into the
admissibility manifold. The listener reconstructs from the manifold back into
semantic space. The gap between the projected trajectory and the reconstructed
one is the fundamental source of communicative failure --- not failure of
information transmission in the Shannon sense, but failure of constraint
alignment: the listener's reconstructed trajectory exits the admissibility
manifold that the speaker was targeting.

\begin{definition}[Conversational repair]
A \textit{conversational repair} is an utterance $r$ whose function is to
return the discourse to the admissibility manifold after a failed transmission.
Repair has the structure of the repair process from the repair ontology:
$r: X \setminus \mathcal{A} \to \mathcal{A}$, mapping an inadmissible
discourse state back into the admissibility region. Repair may target the
speaker's prior utterance (self-initiated self-repair), the listener's
interpretation (other-initiated other-repair), or the conversational
trajectory (metalinguistic repair).
\end{definition}

The repair-rate matching condition from the repair ontology applies directly:
a conversation is stable if and only if the rate of repair $r_R$ meets or
exceeds the rate of semantic drift $r_D$:
\[
r_R(x, t) \geq r_D(x, t) \quad \text{for all } x \in \mathcal{A},\ t \geq 0.
\]
A conversation that exceeds its repair capacity --- in which miscommunications
accumulate faster than repair exchanges can correct them --- will exit the
admissibility manifold. The repair rate is a function of the participants'
shared context, their metalinguistic resources, and the structural accessibility
of repair in the discourse genre.

\begin{theorem}[Admissibility Preservation]
Let $x_0 \in \mathcal{A}$ be an initial discourse state and let
$u_1, u_2, \ldots, u_n$ be a sequence of utterance operators. The discourse
trajectory
\[
x_k = u_k(x_{k-1}), \quad k = 1, \ldots, n,
\]
remains coherent if and only if $x_k \in \mathcal{A}(x_{k-1})$ for every $k$,
equivalently $\rho(x_k) \geq 1$ for every state in the trajectory.

If there exists a smallest index $m$ such that $\rho(x_m) < 1$, then
$x_m \notin \mathcal{A}(x_{m-1})$: the discourse has exited the admissibility
manifold. All subsequent discourse states $x_{m+1}, \ldots, x_n$ either
depend on repair operators $r_k: X \setminus \mathcal{A} \to \mathcal{A}$
restoring the trajectory to $\mathcal{A}$, or they remain outside
$\mathcal{A}$ and the conversation is incoherent from step $m$ onward.
\end{theorem}

\begin{proof}
Discourse coherence at step $k$ is defined as $x_k \in \mathcal{A}(x_{k-1})$.
The trajectory remains coherent at every step iff this condition holds for
all $k \in \{1, \ldots, n\}$, which by the definition of $\rho$ is equivalent
to $\rho(x_k) = S_{\mathrm{support}}(x_k)/\theta(x_k) \geq 1$ for all $k$.
If $\rho(x_m) < 1$ for some $m$, then $x_m \notin \mathcal{A}(x_{m-1})$ by
definition. The utterance $u_m$ has mapped the discourse to an inadmissible
state. For $k > m$, the inherited starting state $x_m$ is outside the
admissibility manifold; subsequent operators either perform repair
($r: X \setminus \mathcal{A} \to \mathcal{A}$, restoring $\rho \geq 1$) or
propagate the inadmissible state, producing further boundary violations.
\end{proof}

\begin{remark}
The Admissibility Preservation Theorem is the discourse-level analogue of
the Semantic Stability Theorem's node-level condition $\rho(v) \geq 1$ and
the branch stability index $\rho_{\min}(\mathcal{B}) \geq 1$. The difference
is temporal: the Semantic Stability Theorem operates over a static tree
structure; the Admissibility Preservation Theorem operates over a dynamic
trajectory through discourse state space. The same repair ratio $\rho$
governs both. The theorem unifies discourse coherence, conversational repair,
CLIO projection, language change (as the evolution of $\mathcal{A}$ across
time), and language death (as the condition $\rho_{\min} < 1$ propagating
through the construction dependency network) under a single formal result.
\end{remark}

\section{Arabic: Root as Attractor, Pattern as Repair Operator}

Arabic root-and-pattern morphology is the most formally transparent
instantiation of the constraint-field approach to grammar, and its
architecture repays careful analysis. The companion paper \textit{Embodied
Constraint and Minimal Assembly} develops a cost-distinctness model of Arabic
script stabilization in which the primitive quantities are production cost
$C$, perceptual distinctness $D$, and system stability $S = f(-C, D)$.
That paper treats the template as a morphological operator $T: R \to W$
mapping a root to a word form. The present section performs the reinterpretation:
\[
T: \mathcal{A}_t \to \mathcal{A}_{t+1},
\]
where the template is a repair operator acting on a discourse admissibility
manifold rather than a combinatorial generator acting on symbolic structures.
The root is not a lexical seed; it is a semantic attractor. The word is not
generated; it is stabilized. This shift from production to maintenance as the
organizing primitive changes what requires explanation and what counts as
explanation.

\subsection{The Standard Account and Its Limitations}

The standard morphological account treats the triliteral root as a semantic
substrate and the pattern (or binyan) as a template that generates word forms
by inserting the root consonants into fixed vowel-consonant frames. The root
\textbf{k-t-b} (associated with writing) generates: \textit{kataba} (he wrote),
\textit{kit\={a}b} (book), \textit{k\={a}tib} (writer, active participle),
\textit{makt\={u}b} (written, passive participle), \textit{maktab} (office,
place of writing), \textit{maktaba} (library), \textit{kit\={a}ba} (writing,
the act), and so on through Forms I--X of the binyan system.

On the generative account, the root is a lexical seed and the patterns are
derivational rules. The morphological system is a productive generator of
word forms. A word is well-formed if it is derivable from an attested root
by an admissible pattern. \textit{Embodied Constraint} formalizes this
as $T = (C_1, C_2, C_3) \mapsto W$, where the template introduces fixed
vocalic insertions, optional consonantal augmentation, and prefixation into
root consonant slots. The generative capacity is constrained:
$|\mathcal{W}| \leq |\mathcal{T}|$, with each template encoding a predictable
semantic modulation.

This account is accurate as a description of the combinatorial productivity
of the system. But it misses the dynamical question: why does the system
have this architecture rather than another? Why does Arabic maintain a stable
consonantal root that survives across phonological and morphological
transformations? What work is the root doing that justifies the cost of
preserving it? The cost-distinctness model answers this in static terms ---
the root is the motor-efficient substrate; dotting provides low-cost
informational refinement --- but does not explain why the system maintains
the root's semantic coherence across derivations over time.

\subsection{The Root as Semantic Attractor}

The constraint-field answer is that the root is a semantic attractor: a high-
density node in the semantic field that exerts repair pressure on its derived
forms, maintaining semantic coherence across contextual deformation.

\begin{definition}[Root attractor and root support]
A triliteral root $R$ defines a \textit{semantic attractor field} $\Phi_R:
\mathcal{L} \to \mathbb{R}_{\geq 0}$ over the lexical space $\mathcal{L}$,
where $\Phi_R(w)$ measures the degree to which word $w$ remains semantically
anchored to $R$'s core meaning field. The \textit{root support} of a derived
form $w$ is:
\[
S_R(w) = \sum_{p \in \mathrm{Path}(R, w)} \alpha_p\, w_p,
\]
where the sum runs over ancestral derivational paths from $R$ to $w$,
$\alpha_p$ is the attenuation along path $p$, and $w_p$ is the admissibility
weight of path $p$. A derived form $w$ is \textit{semantically alive} with
respect to its root when $S_R(w) \geq \theta_w$, where $\theta_w$ is the
minimum root support required for $w$ to participate coherently in the root's
semantic field.
\end{definition}

The patterns (morphological templates) are, on this account, repair operators:
structured deformations of the root that adapt it to specific grammatical
contexts while preserving semantic continuity with the root field. Form I
(\textit{fa\c{c}ala}: basic verb) is the identity repair --- minimal
deformation, maximum root visibility. Form II (\textit{fa\c{c}\c{c}ala}:
intensive/causative) intensifies the root meaning with gemination. The
gemination operator $G(C_i) = C_i C_i$ defined in \textit{Embodied
Constraint} is, in the constraint-field reinterpretation, not just a
syllabic transformation but a semantic intensification repair: by doubling
$C_2$, the pattern increases $\Phi_R$ at the derived form, strengthening
its connection to the root field rather than merely modifying syllabic
architecture. Form X (\textit{istaf\c{c}ala}: estimative/requestive)
introduces the most complex deformation while maintaining the root's semantic
core through the invariant consonantal skeleton $\pi(R) = (C_1, C_2, C_3)$.

\begin{proposition}[Pattern as constraint-preserving deformation]
Each morphological pattern $P$ is a repair operator:
\[
P: \mathcal{A}(R) \to \mathcal{A}(R, G),
\]
mapping the root's semantic admissibility region to the intersection of the
root's semantic field with the grammatical context $G$ (verb, noun, active
participle, place noun, etc.). A pattern is productive when the intersection
$\mathcal{A}(R) \cap \mathcal{A}(G)$ is non-empty: there exist forms that
simultaneously satisfy the root's semantic constraints and the grammatical
context's requirements. The root-order invariance $\pi(R) = (C_1, C_2, C_3)$
proved in \textit{Embodied Constraint} is, in the constraint-field framework,
the formal condition for semantic continuity: the root attractor's identity
is preserved across morphological deformations precisely because the
consonantal skeleton is invariant.
\end{proposition}

The system's resistance to lexical attrition is a direct consequence of the
root-as-attractor structure. The root provides continuous repair pressure on
its derived forms, maintaining semantic coherence against the attenuation that
would otherwise disconnect a word from its semantic ancestry.

In languages without a visible root system, this attenuation runs to
completion unchecked. The general mechanism is: a derived word $w$ undergoes
repeated use in semantic contexts that are progressively distant from its
etymological source. Each such use increases the attenuation $\alpha$ along
the derivational path from the source to $w$, reducing $S_R(w)$ incrementally.
When the source is invisible in the word's form --- because the morphological
connection has been obscured by phonological change, borrowing, or
orthographic opacity --- no repair pressure maintains the semantic connection.
The word's semantic admissibility region drifts from the source field entirely.
The word persists as a stored structure carrying no active connection to its
etymological origin. In Arabic, the consonantal skeleton in the orthography
continually performs a repair function: every occurrence of a word containing
\textbf{k-t-b} activates the root field, maintaining $S_R(w) \geq \theta_w$
for the entire family. The writing system is doing semantic repair work that
English orthography, with its opacity to morphological structure, cannot do.

\subsection{The Incremental Augmentation Principle as Repair Efficiency}

\textit{Embodied Constraint} establishes an Incremental Augmentation Principle
formalized as:
\[
\frac{I(A(X)) - I(X)}{C(A)} \geq \kappa,
\]
where $C(A)$ is the production cost of augmentation $A$, $I$ is informational
distinctness, and $\kappa > 0$ is an efficiency threshold. The principle
states that augmentations persist historically when they deliver high
informational differentiation at low motor cost.

In the constraint-field reinterpretation, this principle is a special case
of the repair efficiency condition: an augmentation persists because it
successfully reduces $S_{\mathrm{RSVP}}$ at the affected nodes --- it
narrows the accessible future volume of interpretation --- without requiring
high repair cost. The dot on \textbf{b} (\textit{ba}) versus undotted
\textbf{t/th} (\textit{ta}/\textit{tha}) achieves $C(A) \approx \epsilon$
while achieving $D(g_1, g_2) \geq D_{\min}$: a minimal repair operation
that prevents the interpretation from exiting the admissibility manifold
into an ambiguous region. The sukūn --- the marker of vowel absence ---
is in this framework an explicit null-state repair: it prevents the default
assumption of vowel continuation from generating inadmissible interpretations.
The shaddah (gemination marker) is a temporal extension repair: a single
graphic mark that modifies syllabic structure across two adjacent syllables
by encoding consonantal doubling, maintaining admissibility at the morphological
boundary without introducing new graphemic primitives.

The repair interpretation explains why the augmentation set $\mathcal{A}$
stabilizes at precisely the inventory it does. The stability criterion
$\partial S / \partial g_i \approx 0$ in \textit{Embodied Constraint} is,
in constraint-field terms, the condition that no further repair operations
are available that would increase the system's admissibility without exceeding
the repair cost threshold. The script has reached a local repair equilibrium:
an attractor state of the repair dynamics in which the system is simultaneously
motor-efficient, perceptually discriminable, and semantically coherent.

\subsection{Dotting as Boundary Preservation}

The diacritical points (nuqat) that distinguish letter pairs sharing a rasm
(consonantal skeleton) --- the \textbf{d/dh} pair sharing one base, the
\textbf{r/z} pair sharing another, the \textbf{s/sh} pair sharing a three-
toothed form --- are not merely phonemic distinguishers. In the constraint-field
framework, they are boundary-preservation operators: graphic marks that prevent
manifold folding in the interpretation space.

\textit{Embodied Constraint} models this as $g_1 = b$, $g_2 = d(b)$ where
$d$ is the dot augmentation satisfying $C(d) \approx \epsilon$ and
$D(g_1, g_2) \geq D_{\min}$. The constraint-field reinterpretation is: the
dot reduces $S_{\mathrm{RSVP}}$ at the graphemic node by collapsing the
accessible future volume of interpretation from a branching ambiguous state
into a single determinate trajectory. Without the dot, the interpretation
space bifurcates at each ambiguous letter; the dot performs a CLIO projection,
collapsing the branching into a single admissible path.

The $\mathbf{r/z}$ pair is especially informative. \textit{Embodied Constraint}
notes that the skeletal form is identical while the dot encodes a manner
contrast involving sustained turbulent airflow (the voiced alveolar fricative
versus the dental stop). In constraint-field terms, the dot marks the boundary
between two distinct regions of the admissibility manifold: the stop-manner
region and the fricative-manner region. Without the dot, an utterance
containing this letter occupies a point equidistant from both boundary regions
--- a maximally unstable repair position, requiring the listener to perform
external disambiguation from context. The dot is the writer's repair operation,
performed in advance of transmission, that moves the interpretation away from
the boundary and into the interior of the correct manifold region.

The non-connecting letters (\textit{alif, dal, dhal, ra, zayn, waw}) function
as semantic fault barriers: strategic discontinuities in the script that
arrest ambiguity propagation. \textit{Embodied Constraint} models these as
constraint insertions that preserve minimum perceptual distance
$D_{\mathrm{join}} \geq D_{\min}$ by imposing $J(g_i, g_{i+1}) = 0$ for
specific $g_i$. The constraint-field reinterpretation: these letters are
points in the word where the admissibility manifold would fold if joining
continued uninterrupted --- where adjacent rasm forms, joined, would collapse
into a single interpretive region that the listener cannot disambiguate.
The non-connecting letter is not a motor convenience; it is a repair
intervention built into the orthographic infrastructure, preventing the
interpretation from crossing an admissibility boundary by forced segmentation.

The positional allograph system --- the reduction from context-sensitive
ligatures ($\mathcal{O}(n^2)$ possible forms) to positional forms
($\mathcal{O}(4n)$) --- is in \textit{Embodied Constraint} a dimensional
reduction that decreases glyph-form entropy $H$ while preserving cursive
continuity. In the constraint-field framework, this is the system achieving
a stable repair equilibrium: the high-entropy ligature system required too
much disambiguation repair from readers and too much production precision
from writers; the positional system achieves the same admissibility guarantees
at lower repair cost. The compression from $H_{\mathrm{context}}$ to
$H_{\mathrm{positional}}$ is not a loss but an efficiency gain in the repair
dynamics.

\subsection{Lexicalization as Repair Collapse}

The constraint-field framework gives a precise account of lexicalization ---
the process by which a derived form becomes semantically independent of its
root --- as a special case of repair collapse.

When a derived form $w$ undergoes repeated use in contexts that are
semantically distant from the root field $\Phi_R$, the attenuation $\alpha$
along the path from $R$ to $w$ increases. The root support $S_R(w)$ declines.
When $S_R(w)$ crosses the threshold $\theta_w$, the form exits the root's
semantic admissibility region:
\[
S_R(w) < \theta_w \implies w \notin \mathcal{A}(R).
\]
At this point $w$ is semantically dead with respect to the root: it continues
to exist as a lexical item but no longer participates in the root's repair
network. It has become an isolated stored structure, carrying the consonantal
form of the root without the semantic support that the root field provides.

This is the formal mechanism of lexicalization. The Arabic word
\textit{waz{\=\i}r} (minister, vizier) derives from the root \textbf{w-z-r}
(to bear a burden), but in ordinary use it no longer activates the burden-
bearing semantic field. The root support has fallen below threshold. The word
persists as a stored structure; the root attractor no longer repairs the
semantic connection.

The Arabic system's resistance to lexicalization --- relative to non-root-
based languages --- is a direct consequence of the orthographic repair
infrastructure. The consonantal skeleton prevents the attenuation from
completing: as long as the root is visible in the word's written form, the
root attractor maintains some repair pressure, slowing the descent of
$S_R(w)$ below $\theta_w$. This is the constraint-field interpretation of
\textit{Embodied Constraint}'s central thesis: the Arabic script's layered
architecture (skeletal forms, dotting, diacritics, positional regularization)
stabilizes as a constrained symbolic equilibrium because each layer is a
repair operation that maintains the admissibility of the next layer's output.
The script does not merely encode language; it continuously repairs the
semantic coherence of the language it encodes.

\section{Spanish: Multi-Field Coherence}

Spanish verbal morphology illustrates a different mode of the same constraint-
field structure: not the root-attractor system of Arabic but the multi-field
coherence condition that arises when multiple independent constraint dimensions
must be simultaneously satisfied.

\subsection{Agreement as Field Intersection}

The Spanish verbal agreement system requires the verb to be simultaneously
consistent with person, number, tense, aspect, mood, and (in some constructions)
gender. Traditional grammar describes this as rule application: the verb
selects the form that satisfies all the relevant features of the subject.
The constraint-field approach describes it as field intersection: the verb
must occupy a region of the discourse state space that lies within the
admissibility regions simultaneously imposed by each active constraint dimension.

The form \textit{hablo} (I speak, present indicative) is not the result of
applying a set of rules to the verb \textit{hablar} and the features [1st
person, singular, present, indicative]. It is the unique attractor in the
intersection of the following constraint fields:

\begin{align*}
\mathcal{A}(\text{1st person}) &: \text{forms marked for speaker reference} \\
\mathcal{A}(\text{singular}) &: \text{forms unmarked for plurality} \\
\mathcal{A}(\text{present}) &: \text{forms occupying the present tense basin} \\
\mathcal{A}(\text{indicative}) &: \text{forms within the realis mood region}
\end{align*}

The correct form is the one that lies in $\bigcap_i \mathcal{A}(d_i)$, where
$d_i$ ranges over active constraint dimensions. An agreement error is a form
that lies outside this intersection: it satisfies some constraint dimensions
but not all, and the resulting discourse state exits the admissibility manifold.

\subsection{Native Competence as Field Inhabitation}

The constraint-field account explains an otherwise puzzling asymmetry between
native and non-native speaker behavior. Native speakers of Spanish report that
correct agreement \textit{feels} right and incorrect agreement \textit{feels}
wrong --- they do not consult a rule, they sense a violation. Non-native
speakers, by contrast, often describe themselves as checking rules.

On the generative account this is a performance/competence distinction: native
speakers have internalized the grammar and apply it automatically. The constraint-
field account gives a different explanation: native speakers inhabit the discourse
admissibility manifold. They perceive agreement errors as boundary crossings ---
as states outside the admissibility region --- because they have calibrated their
perceptual sensitivity to the field's geometry through years of exposure. They
are not applying rules; they are detecting field violations.

Non-native speakers who report consulting rules are operating outside the field:
they know the grammar's constraints as explicit propositions but have not yet
developed the perceptual sensitivity to the field's geometry. The developmental
trajectory of second-language acquisition is, on this account, the trajectory
of field inhabitation: the learner progressively calibrates their sensitivity
to the constraint field until explicit rule-consultation becomes unnecessary.

\begin{proposition}[Competence as field sensitivity]
Linguistic competence in language $L$ is the capacity to detect whether a
discourse trajectory lies within $\mathcal{A}_L$, the admissibility manifold
of $L$, without explicit computation of the constraints defining $\mathcal{A}_L$.
This capacity is acquired through exposure to the distribution of admissible
and inadmissible trajectories in $L$, which calibrates the learner's
sensitivity to the field geometry.
\end{proposition}

\subsection{Subjunctive as Admissibility Shift}

The Spanish subjunctive mood is particularly illuminating because it cannot
be described as a rule without appealing to semantic and pragmatic conditions
that exceed the scope of formal syntax. The subjunctive is required in
embedded clauses under verbs of desire, doubt, denial, and emotion --- but
not under verbs of assertion or perception. The syntactic distribution of the
subjunctive is not formally characterizable without reference to the epistemic
and affective commitments of the main clause.

The constraint-field account handles this naturally. The indicative occupies
the realis region of the discourse admissibility manifold: the region in which
the embedded proposition is asserted as obtaining in the actual or projected
world. The subjunctive occupies the irrealis region: the region in which the
embedded proposition is presented as desired, doubted, denied, or emotionally
evaluated rather than asserted.

Verbs of desire, doubt, and denial require the subjunctive because they shift
the embedded clause into the irrealis region of the admissibility manifold.
The main clause verb is a field operator that transforms the constraint
structure of the embedded clause: \textit{quiero que vengas} (I want you to
come) applies the desire operator, shifting the embedded clause from the realis
to the irrealis region. Using the indicative (\textit{quiero que vienes}) is
a boundary crossing: it places the embedded proposition in the realis region,
which is incompatible with the desire frame established by the main clause.

\section{Conversation as Real-Time Boundary Maintenance}

\subsection{Turn-Taking as Constraint Relay}

A conversation is a sequence of CLIO projections and repair operations in
which participants alternate in applying utterances as operators on the shared
discourse state space. Turn-taking is not merely a coordination mechanism for
orderly speaker alternation; it is the structure through which constraint relay
is maintained.

Each speaker turn is a commitment: the speaker applies an utterance operator
$u$ that commits the discourse to a new state $u(x)$. This state must lie
within $\mathcal{A}(x)$ to be receivable. But it also establishes the
constraint structure for the next speaker's turn: the next speaker inherits
$u(x)$ as their starting state and must produce an utterance that maintains
coherence from $u(x)$ rather than from $x$.

Turn-taking failures --- interruptions, topic hijacks, non-sequiturs ---
are not merely violations of conversational norms. They are constraint
structure violations: they apply utterance operators to discourse states
other than the inherited state, producing trajectories that are inadmissible
given the established discourse context. The conversational repair that follows
such failures (the interrupted speaker resuming, the hijacked topic being
reinstated, the non-sequitur being flagged) is the conversation's repair
process at work: returning the discourse trajectory to the admissibility
manifold.

\subsection{Genre as Admissibility Architecture}

Different discourse genres --- casual conversation, academic lecture, legal
testimony, religious ritual, intimate disclosure --- maintain different
admissibility manifolds. The same utterance may be admissible in one genre
and inadmissible in another, not because the sentence is syntactically
different but because the genre establishes different constraint structures.

Genre is therefore an admissibility architecture: a persistent constraint
structure that shapes the discourse manifold for a class of communicative
events. Participants entering a genre do not merely follow its conventions;
they enter its admissibility manifold and calibrate their utterance operators
to the field's geometry. A participant who produces a genre-inappropriate
utterance has not violated a convention; they have crossed a boundary. The
conversational repair that follows (reframing, correction, discomfort) is
the genre's repair infrastructure responding to the boundary crossing.

\section{Language Change as Manifold Deformation}

\subsection{Constructions and Their Repair Ratios}

Applying the Semantic Stability framework to grammatical constructions yields
a precise account of language change as manifold deformation.

\begin{definition}[Constructional repair ratio]
For a grammatical construction $c$, the \textit{constructional repair ratio} is:
\[
\rho(c) = \frac{S(c)}{\theta(c)},
\]
where $S(c)$ is the \textit{communicative support} of $c$ --- the degree to
which the construction is actively reinforced through successful use in
communicative exchanges --- and $\theta(c)$ is the \textit{maintenance
threshold} of $c$ --- the minimum support required for $c$ to remain within
the admissibility manifold of the language.
\end{definition}

A construction is grammatically alive when $\rho(c) \geq 1$. It is undergoing
attrition when $\rho(c) < 1$: the communicative support has fallen below the
maintenance threshold, and the construction is in the process of exiting the
language's admissibility manifold.

Language change, on this account, is the evolution of the admissibility manifold
through the differential attrition and emergence of constructions. A construction
that consistently produces admissible discourse states accumulates repair support
and deepens its basin in the manifold. A construction that fails to produce
admissible states loses support and approaches the boundary. When $\rho(c)$
crosses $1$, the construction undergoes the same failure-front propagation
established for semantic trees: forms that depended on the failing construction
for their own admissibility lose support, and the collapse propagates through
the grammatical dependency structure.

\subsection{Sound Change as Repair of Articulatory Effort}

Sound change is the clearest case of manifold deformation as repair. The
pressures driving sound change are well-understood: assimilation (adjacent
sounds influence each other), lenition (intervocalic consonants weaken),
merger (phonemically distinct sounds become indistinguishable in certain
contexts), and chain shifts (systematic reorganizations of the vowel or
consonant space). These are all repair operations applied to the articulatory-
perceptual constraint manifold.

Assimilation is repair of articulatory discontinuity: the effort required to
produce a sequence of sounds with very different articulations is reduced by
partially assimilating adjacent sounds to each other. The result is a discourse
trajectory that remains within the articulatory admissibility manifold with
lower articulatory cost.

Lenition is repair of constraint excess: intervocalic consonants carry minimal
distinguishing load in many phonological contexts, and the effort required to
maintain their full stop or fricative articulation exceeds the constraint
maintenance benefit. The lenited form costs less effort while remaining within
the perceptual admissibility manifold of the listener.

Grammaticalization is repair of semantic underspecification: a construction
that repeatedly appears in contexts where its meaning is semantically bleached
--- used for a grammatical function rather than its lexical meaning ---
accumulates the grammatical reading as an admissible trajectory and loses the
lexical reading as the context consistently selects against it. The
construction's admissibility region shifts from the lexical to the grammatical
manifold.

\section{Language Death as Repair Collapse}

The Semantic Stability Theorem, applied to language as a whole, gives the
following account of language death.

A language is a repair equilibrium: a constraint-field system in which the
admissibility manifold $\mathcal{A}_L$ is maintained by the continuous repair
activity of the speech community. Every successful communicative exchange
is a repair operation that reinforces the constructions it employs. Every
failed exchange that is subsequently repaired strengthens the repair
infrastructure. The language persists as long as the repair rate meets or
exceeds the degradation rate.

Language death is repair collapse: the failure of the repair rate to meet
the degradation rate, producing a failure-front propagation from peripheral
constructions toward the core vocabulary and grammar.

The branch stability index $\rho_{\min}(\mathcal{L})$ --- the minimum repair
ratio across all constructions in the language --- is the survival condition.
When $\rho_{\min}(\mathcal{L}) < 1$, at least one construction has crossed
the admissibility boundary, and the failure-front dynamics of the Semantic
Stability Theorem apply.

\begin{corollary}[Language death as repair failure]
A language undergoes irreversible collapse when $\rho_{\min}(\mathcal{L}) < 1$
for a connected subgraph of the constructional dependency network. The
collapse initiates at the peripheral constructions (those with the smallest
repair surplus $S(c) - \theta(c)$), propagates inward through the dependency
structure, and terminates when either external support is introduced or
the core vocabulary falls below its admissibility threshold.
\end{corollary}

The Critical Repair Divergence theorem applies here with particular force. As
a language approaches the admissibility boundary --- as $\rho_{\min}$ approaches
$1$ from above --- the repair intensity required to maintain the remaining
constructions diverges. Late-stage language revitalization is extremely costly
not because the language is inherently difficult to revitalize but because the
system is operating near its admissibility boundary, where $\iota \to \infty$.
The cost of maintaining each surviving construction approaches infinity as the
surrounding repair network collapses.

This also explains why revitalization requires reconstructing the repair
infrastructure rather than merely reintroducing the content. A dictionary of
a dying language contains the nodes $\ell(v)$ but not the constraint
infrastructure $S_R(w) \geq \theta_w$ that makes each node semantically alive.
A child who learns vocabulary from a dictionary but does not acquire the
discourse manifold is not a speaker of the language; they are a user of stored
structures disconnected from the living repair equilibrium.

\section{RSVP as Field-Theoretic Synthesis}

The discourse manifold, the utterance operator, the repair ratio, and the
Admissibility Preservation Theorem are sufficient to account for the phenomena
examined in the preceding sections. The Relativistic Scalar-Vector Plenum
framework does not introduce new empirical claims about language; it provides
a field-theoretic substrate in which the discourse-level quantities acquire
continuous geometric descriptions. RSVP appears here as a consequence of the
framework established through the case studies, not as a starting assumption.

\begin{definition}[RSVP-discourse correspondence]
The RSVP field triple $(\Phi, v, S_{\mathrm{RSVP}})$ is identified with the
discourse field by the following explicit equalities, forced by the requirement
that the RSVP stability condition coincide with the Admissibility Preservation
Theorem:
\begin{align}
\Phi(x) &= S_{\mathrm{support}}(x)
  \quad \text{(semantic density = support reaching state } x\text{)} \\
|v_{ij}| &= \alpha(i,j)
  \quad \text{(vector magnitude = constraint transmission coefficient)} \\
S_{\mathrm{RSVP}}(x) &= -\log \rho(x)
  \quad \text{(entropy field = log-inverse repair ratio)}
\end{align}
\end{definition}

The identification $S_{\mathrm{RSVP}} = -\log\rho$ is the structurally
significant one. It has three immediately verifiable properties: when
$\rho(x) > 1$ (the discourse state is inside the admissibility manifold),
$S_{\mathrm{RSVP}} < 0$ --- the entropy is negative, encoding the constraint
surplus that keeps the state interior. When $\rho(x) = 1$ (exactly at the
boundary), $S_{\mathrm{RSVP}} = 0$. When $\rho(x) < 1$ (outside the
manifold), $S_{\mathrm{RSVP}} > 0$ --- positive entropy signals the expansion
of interpretation volume beyond what the constraint structure can contain.
The Admissibility Preservation Theorem translates directly: discourse coherence
is maintained iff $S_{\mathrm{RSVP}}(x_k) \leq 0$ for all steps $k$.

\begin{proposition}[RSVP admissibility condition]
Under the identification above, an utterance $u$ is admissible in discourse
state $x$ if and only if $S_{\mathrm{RSVP}}(u(x)) \leq 0$.
\end{proposition}

\begin{proof}
$S_{\mathrm{RSVP}}(u(x)) = -\log\rho(u(x)) \leq 0$ iff $\rho(u(x)) \geq 1$
iff $u(x) \in \mathcal{A}(x)$.
\end{proof}

The scalar field $\Phi = S_{\mathrm{support}}$ provides the field geometry of
the discourse. High $\Phi$ at a discourse state means richly constrained
interpretation: a topic-initial position, a loaded term, a well-established
presupposition. Low $\Phi$ means loosely constrained interpretation: an
open-ended question, a genre-boundary crossing, an unfamiliar register.
The vector field magnitude $|v_{ij}| = \alpha(i,j)$ is the constraint
transmission coefficient between adjacent discourse states: how much of the
constraint established at $i$ reaches $j$. Low $|v|$ on an edge means
high attenuation --- culturally distant participants, mismatched registers,
or the communicative equivalent of a non-connecting letter.

\subsection{Orthographic Conventions as Generalized Repair Infrastructure}

The Arabic analysis reveals a generalization that extends to all writing
systems and to many non-orthographic symbolic conventions. The structural
features that the Arabic system uses to maintain its admissibility manifold
--- dots, non-connecting letters, positional regularization, diacritic
layers --- all function as repair infrastructure built into the medium.

Once this is seen, the same structure becomes visible elsewhere. Word spacing
in most alphabetic scripts is a non-connecting operator: a mandatory
segmentation that prevents ambiguity from propagating across word boundaries,
exactly as the non-connecting letters in Arabic prevent ambiguity from
propagating within words. Without word spaces, a Latin-script text becomes
interpretively ambiguous at every boundary: \textit{therapist} versus
\textit{the rapist}, \textit{nowhere} versus \textit{now here}. The space is
not a convention; it is a boundary-preservation operator that maintains
$D_{\mathrm{join}} \geq D_{\min}$ at word boundaries.

Punctuation is turn-taking infrastructure for written discourse: the period
terminates a clause bubble, the comma maintains syntactic connection across
a pause, the colon announces a constraint-commitment (what follows will be
constrained by what preceded), the quotation mark shifts the admissibility
manifold to a nested register. Capitalization distinguishes the start of a
new repair equilibrium (sentence-initial, proper name) from continuation
within an existing one.

Paragraph breaks, section headings, and chapter divisions are higher-order
fault barriers: points at which the admissibility manifold resets to a lower
constraint density, allowing readers to re-enter the document's semantic
field without carrying all prior constraints. A text without paragraph
breaks requires readers to maintain the entire prior constraint structure
in working memory; the break is a repair operation that reduces the cognitive
cost of re-entering the manifold.

In the RSVP field interpretation: every orthographic convention that reduces
ambiguity is a point intervention that lowers $S_{\mathrm{RSVP}}$ at the
affected node. Every convention that marks a boundary (word space, period,
paragraph break) is a CLIO projection that collapses the branching
interpretation space into a determinate trajectory. Writing systems are not
representations of speech; they are layered repair infrastructures for
discourse admissibility, and their conventions are the materialized residue
of accumulated repair operations that have been found efficient enough to
warrant standardization.

\begin{remark}[Companion paper]
The analysis of writing systems as boundary-preservation technologies is
sufficiently rich to constitute a separate inquiry. A companion essay ---
\textit{Repair Infrastructures: Writing Systems as Boundary-Preservation
Technologies} --- would extend the Arabic analysis to Latin script conventions,
to ideographic and logographic systems, to the history of punctuation, and
to digital typography. The present section establishes the theoretical
connection; the companion would develop the empirical breadth.
\end{remark}

\section{Gesture Before Symbol: The Root Level}

The gesture-before-symbol thesis --- that meaning is prior to its linguistic
encoding, and that symbols are repair layers applied to pre-symbolic meaning
--- specifies the position of the discourse manifold in the overall hierarchy
of semantic structure.

\begin{definition}[Semantic hierarchy]
The semantic hierarchy from meaning to utterance has the following structure:
\begin{align*}
\text{Gesture} &: \text{pre-symbolic meaning grounded in shared attention and action} \\
\downarrow & \\
\text{Constraint field} &: \text{the admissibility structure that makes utterances receivable} \\
\downarrow & \\
\text{Root attractor} &: \text{semantic anchor maintaining coherence across derivations} \\
\downarrow & \\
\text{Morphological repair} &: \text{pattern operators adapting the root to grammatical context} \\
\downarrow & \\
\text{Utterance} &: \text{the committed trajectory within the admissibility manifold}
\end{align*}
\end{definition}

Each level of the hierarchy is a repair layer applied to the level above it.
Grammar does not generate meaning from rules; grammar repairs the transmission
of pre-symbolic meaning into the social medium of language, maintaining
semantic coherence against the attenuation introduced by the arbitrariness
of the symbol, the distance between interlocutors, and the variation in
individual context.

This hierarchy inverts the generative direction: rather than meaning being
produced by the grammar, meaning is the source, and the grammar is the
infrastructure that maintains access to it. A language is not a machine for
producing sentences. A language is a repair ecology that preserves shared
access to the pre-symbolic meaning that speakers bring to the conversation.

The consequence for language acquisition is significant. A child acquiring
language is not learning to generate sentences; they are calibrating their
sensitivity to the discourse admissibility manifold. The critical period for
language acquisition is not the period during which grammar rules are
internalizable but the period during which the field geometry can be calibrated
at the perceptual level. A child who acquires a language during the critical
period inhabits its admissibility manifold; a child who acquires it afterward
must compensate with explicit rule knowledge, which is a less efficient
repair infrastructure.

\section{Conclusion}

The constraint-field theory of language proposed in this essay makes three
central claims that distinguish it from the generative tradition.

The first is that the primitive object is the discourse admissibility manifold,
not the sentence or the derivation. Grammar is the constraint structure that
defines the manifold; syntax is the record of how those constraints propagate
through utterances; grammaticality is a property of trajectories rather than
objects.

The second is that linguistic competence is field sensitivity, not rule
knowledge. Native speakers inhabit the admissibility manifold; they detect
boundary crossings as violations rather than computing rule failures. The
developmental trajectory of acquisition is the trajectory of field inhabitation.
The asymmetry between native and non-native speaker behavior is the asymmetry
between manifold inhabitation and explicit constraint computation.

The third is that language change is manifold deformation driven by the
differential repair rates of grammatical constructions. Sound change, semantic
change, and grammaticalization are all repair operations on the articulatory,
semantic, and grammatical constraint manifolds respectively. Language death
is repair collapse. Language revitalization requires reconstructing the repair
infrastructure, not reintroducing the content.

The Arabic root-and-pattern system is the most formally transparent
instantiation of this framework: the root as semantic attractor providing
continuous repair pressure on its derived forms, the patterns as repair
operators adapting the root to grammatical context, the orthography as
a repair infrastructure that prevents lexicalization by keeping the root
visible at every use. Spanish verbal agreement is an instance of multi-field
coherence: the correct form is the unique attractor in the intersection of
simultaneously active constraint fields, and native competence is the
perceptual sensitivity to that intersection's geometry.

Grammar does not generate utterances. Grammar maintains the admissibility
structure that makes utterances receivable. The utterance is not the output
of a production system; it is the committed trajectory of a repair process
that seeks to transmit meaning through a medium --- the social-linguistic
field --- that continuously introduces attenuation. Every successful
utterance is a repair. Every successful conversation is a shared repair
equilibrium, maintained in real time by participants who have calibrated
their sensitivity to the same constraint field.

The river is not the water. The grammar is not the sentence. The grammar
is the work that makes the sentence possible.

\bigskip

\noindent\rule{0.4\textwidth}{0.4pt}

\vspace{0.5em}
\noindent{\small Flyxion / Independent Researcher. 2026.}

\end{document}

